\section{Metodología y herramientas}
\label{sec:metodologia}

\paragraph{Cadena de validación.} Cada etapa se contrasta con la anterior, de modo que toda discrepancia pueda atribuirse a una causa
concreta:
\begin{enumerate}
  \item \textbf{Función objetivo} $H_t(s)$: aproximación obtenida con \texttt{scipy.signal} en forma ceros--polos--ganancia (zpk),
        validada contra las expresiones analíticas de cada aproximación.
  \item \textbf{Función del circuito} $H_c(s)$: análisis nodal simbólico con SymPy~\cite{meurer2017} y AO ideal; se identifica
        $\omega_0$, $Q$ y ganancia de cada sección y se compara con $H_t(s)$.
  \item \textbf{Componentes comerciales}: recálculo de $H_c(s)$ con valores normalizados (series E) y estudio de sensibilidades.
  \item \textbf{LTspice}~\cite{ltspice}: circuito con AO ideal y con el modelo del fabricante; análisis AC y transitorio.
  \item \textbf{Prototipo}: medición de amplitud y fase en las bandas de paso, de transición y de rechazo.
\end{enumerate}

\paragraph{Herramientas.} Python con NumPy~\cite{harris2020}, SciPy~\cite{virtanen2020}, SymPy~\cite{meurer2017},
Matplotlib~\cite{hunter2007} y pandas (notebooks ejecutables en Google Colab), LTspice para la simulación circuital, \LaTeX{} para el
informe y un repositorio Git con las fuentes, los datos y un \texttt{README} con instrucciones de reproducción y versiones de las
bibliotecas. Las figuras y tablas numéricas de este informe se generan directamente desde los notebooks (carpetas
\texttt{figuras/} y \texttt{resultados/}), de modo que tablas, figuras y valores citados en el texto salen de los mismos cálculos y no de transcripciones manuales.

\paragraph{Criterios numéricos.} Se trabaja en forma zpk y en secciones de segundo orden en lugar de coeficientes polinomiales, que
están mal condicionados para órdenes altos. El retardo de grupo se calcula analíticamente a partir de polos y ceros
(\cref{eq:retardo}), con $\omega$ en rad/s, y se contrasta con la derivada numérica de la fase desenrollada.
